\rtmtight
\begin{tblr}{width=\textwidth,colspec={|Q[c,m,2.1cm]|X[l,m]|},hlines,vlines,hspan=minimal,rowsep=1.5pt,colsep=3pt,row{1}={bg=headgray,font=\bfseries}}
\SetCell{c,m}\hfil\logo\hfil & \textbf{POLITEKNIK NEGERI MALANG}\par\textbf{JURUSAN TEKNOLOGI INFORMASI}\par\textbf{PROGRAM STUDI: D4 TEKNIK INFORMATIKA} \\
\SetCell[c=2]{l} \textbf{RENCANA TUGAS MAHASISWA (RTM) DAN RUBRIK PENILAIAN} & \\
Nama Mata Kuliah & Pengolahan Citra dan Visi Komputer \\
Kode Mata Kuliah & RTI255007 \quad \textbf{sks / Jam perminggu:} 3 SKS/6 jam \quad \textbf{Semester:} 5 \\
Dosen Pengampu & Tim Pengajar Program Studi D4 TI \\
\end{tblr}
\assessmentblock{Tugas Pengolahan Citra Dasar}{Case Method dan praktik implementasi}{SCPMK0705-04301, SCPMK1004-04303, SCPMK0706-04302}{Mahasiswa menyelesaikan tugas mengimplementasikan teknik pengolahan citra menggunakan OpenCV, mengevaluasi kualitas citra, dan membangun implementasi CNN awal.}{Menganalisis kebutuhan, mengimplementasikan teknik citra dengan OpenCV, mengukur metrik kualitas, mengimplementasikan CNN awal, dan mendokumentasikan.}{Kode Python/OpenCV, laporan evaluasi kualitas citra, implementasi CNN, dan/atau bahan presentasi.}{Implementasi teknik filtering dan segmentasi & 12 & Teknik filtering (Gaussian, Sobel, Canny), segmentasi (thresholding, watershed), dan transformasi diimplementasikan dengan benar menggunakan OpenCV dan hasilnya terdokumentasi. & Implementasi sebagian besar benar tetapi beberapa teknik belum optimal atau hasilnya belum terdokumentasi lengkap. & Teknik pengolahan citra tidak diimplementasikan dengan benar atau kode tidak berjalan. \\ \\
Evaluasi kualitas citra menggunakan metrik kuantitatif & 8 & Metrik kualitas citra (PSNR, SSIM, histogram) digunakan dengan benar, hasilnya diinterpretasikan secara akurat, dan perbandingan antar teknik dilakukan secara sistematis. & Metrik digunakan tetapi interpretasi atau perbandingan antar teknik masih belum lengkap. & Metrik tidak digunakan atau interpretasi tidak mencerminkan pemahaman kualitas citra. \\ \\
Implementasi CNN awal untuk klasifikasi citra & 8 & Arsitektur CNN dasar diimplementasikan dengan benar, proses pelatihan berjalan, dan hasil klasifikasi awal terdokumentasi. & Implementasi CNN ada tetapi arsitektur atau proses pelatihan masih ada kekurangan. & CNN tidak diimplementasikan atau tidak dapat menghasilkan hasil klasifikasi. \\}

\assessmentblock{Kuis 1 Pengolahan Citra Dasar}{Kuis}{SCPMK0705-04301}{Kuis tertulis untuk mengukur pemahaman teknik pengolahan citra dasar menggunakan OpenCV.}{Menjawab soal pilihan ganda dan/atau uraian terkait representasi citra, operasi pixel, filtering, dan deteksi tepi.}{Jawaban tertulis kuis.}{Pemahaman teknik pengolahan citra dasar & 3 & Konsep representasi citra, operasi pixel, histogram, filtering, dan deteksi tepi dijelaskan dengan benar dan contoh penerapannya tepat. & Sebagian besar konsep dipahami tetapi penerapan atau contoh masih kurang lengkap. & Konsep tidak dipahami atau contoh penerapan keliru. \\ \\
Analisis operasi pixel dan histogram & 2 & Operasi pixel (brightness, contrast) dan manipulasi histogram dianalisis dengan tepat pada kasus citra yang diberikan. & Analisis sebagian besar benar tetapi beberapa operasi atau hasil histogram masih keliru. & Analisis tidak tepat atau tidak menunjukkan pemahaman operasi citra. \\}

\assessmentblock{UTS Pengolahan Citra dan Evaluasi Kualitas}{UTS berbasis praktik}{SCPMK0705-04301, SCPMK1004-04303}{Evaluasi tengah semester untuk mengukur kemampuan mengimplementasikan teknik pengolahan citra dan mengevaluasi kualitas hasil pemrosesan visual.}{Mengimplementasikan teknik pengolahan citra pada dataset yang diberikan dan mengevaluasi kualitas hasilnya.}{Kode Python/OpenCV yang berjalan, hasil pemrosesan citra, dan laporan evaluasi kualitas.}{Implementasi teknik pengolahan citra & 8 & Teknik filtering, segmentasi, dan transformasi diimplementasikan dengan benar pada dataset UTS, kode berjalan tanpa error, dan hasilnya sesuai ekspektasi. & Implementasi sebagian besar benar tetapi beberapa teknik menghasilkan output yang tidak optimal atau ada error minor. & Implementasi tidak berjalan atau hasilnya tidak mencerminkan penerapan teknik yang benar. \\ \\
Evaluasi kualitas hasil pemrosesan visual & 7 & Kualitas hasil pemrosesan dievaluasi dengan metrik yang tepat (PSNR, SSIM), diinterpretasikan secara akurat, dan dibandingkan antar teknik secara sistematis. & Evaluasi dilakukan tetapi pemilihan metrik atau interpretasi hasil masih belum lengkap. & Evaluasi tidak dilakukan atau tidak menggunakan metrik yang relevan. \\}

\assessmentblock{Kuis 2 Evaluasi Model Visi Komputer}{Kuis}{SCPMK1005-04304}{Kuis tertulis untuk mengukur pemahaman metrik evaluasi model visi komputer.}{Menjawab soal terkait metrik mAP, IoU, accuracy, precision, recall, dan interpretasinya.}{Jawaban tertulis kuis.}{Pemahaman metrik evaluasi model visi & 3 & Metrik mAP, IoU, accuracy, precision, dan recall dijelaskan dengan benar dan aplikasinya pada kasus deteksi/segmentasi tepat. & Sebagian besar metrik dipahami tetapi penerapan pada kasus spesifik masih kurang tepat. & Metrik tidak dipahami atau penerapannya keliru. \\ \\
Kemampuan analisis kinerja model & 2 & Kinerja model dievaluasi secara komprehensif berdasarkan metrik yang sesuai konteks aplikasi visi komputer. & Analisis kinerja cukup baik tetapi pemilihan metrik atau interpretasi hasilnya belum sepenuhnya tepat. & Analisis tidak tepat atau tidak menggunakan metrik yang relevan. \\}

\assessmentblock{PBL Implementasi Sistem Visi Komputer}{Project Based Learning}{SCPMK0706-04302, SCPMK1005-04304}{Mahasiswa membangun sistem visi komputer berbasis CNN untuk klasifikasi, deteksi objek, atau segmentasi, dan mengevaluasi kinerjanya menggunakan metrik kuantitatif.}{Menganalisis masalah, merancang arsitektur CNN, melatih model, mengevaluasi kinerja, dan mempresentasikan sistem.}{Kode Python/PyTorch, model terlatih, laporan evaluasi, dan bahan presentasi.}{Implementasi sistem visi komputer berbasis CNN & 10 & Sistem CNN diimplementasikan dengan arsitektur yang sesuai (klasifikasi/deteksi/segmentasi), proses pelatihan terdokumentasi, dan model menghasilkan output yang benar pada data uji. & Implementasi CNN berjalan tetapi arsitektur atau proses pelatihan masih ada kekurangan yang mempengaruhi kualitas output. & CNN tidak diimplementasikan dengan benar atau tidak menghasilkan output yang bermakna. \\ \\
Evaluasi kinerja menggunakan metrik kuantitatif & 8 & Kinerja model dievaluasi dengan metrik yang tepat (mAP, IoU, accuracy), hasil diinterpretasikan secara akurat, dan perbandingan dengan baseline dilakukan. & Evaluasi dilakukan dengan metrik yang benar tetapi interpretasi atau perbandingan dengan baseline masih belum lengkap. & Evaluasi tidak dilakukan atau menggunakan metrik yang tidak sesuai konteks. \\ \\
Dokumentasi dan presentasi sistem & 7 & Sistem didokumentasikan lengkap (arsitektur, data, pelatihan, evaluasi) dan dipresentasikan dengan penjelasan yang meyakinkan tentang pilihan teknis. & Dokumentasi dan presentasi cukup baik tetapi beberapa aspek teknis belum dijelaskan dengan lengkap. & Dokumentasi tidak lengkap atau presentasi tidak dapat menjelaskan sistem yang dibangun. \\}

\assessmentblock{UAS Presentasi dan Evaluasi Sistem Visi Komputer}{UAS presentasi dan defense}{SCPMK0706-04302, SCPMK1005-04304}{Mahasiswa mempresentasikan sistem visi komputer yang dibangun, mendemonstrasikan fungsionalitas, dan mempertahankan keputusan teknis berdasarkan evaluasi kinerja.}{Melakukan demo sistem, mempresentasikan hasil evaluasi, dan menjawab pertanyaan teknis.}{Sistem visi komputer final, laporan evaluasi komprehensif, dan bahan presentasi.}{Kelengkapan implementasi sistem visi komputer & 9 & Sistem visi komputer berjalan lengkap dengan pipeline yang benar (preprocessing, model, evaluasi), semua komponen terintegrasi, dan demo berjalan tanpa error. & Sebagian besar komponen berfungsi tetapi beberapa integrasi pipeline atau demo masih ada kekurangan. & Sistem tidak berjalan atau banyak komponen yang tidak terintegrasi. \\ \\
Kualitas evaluasi kinerja model & 8 & Kinerja model dievaluasi secara komprehensif dengan metrik yang tepat, hasilnya dianalisis mendalam, dan keterbatasan model diidentifikasi dengan jelas. & Evaluasi kinerja dilakukan tetapi kedalaman analisis atau identifikasi keterbatasan masih belum lengkap. & Evaluasi tidak komprehensif atau tidak menggunakan metrik yang sesuai. \\ \\
Analisis dan presentasi hasil & 5 & Hasil sistem dijelaskan dengan tepat, keputusan teknis dipertahankan dengan argumen berbasis data, dan pertanyaan dijawab secara meyakinkan. & Presentasi cukup baik tetapi beberapa keputusan teknis belum dapat dipertahankan dengan argumen yang kuat. & Tidak dapat menjelaskan sistem atau menjawab pertanyaan teknis dengan tepat. \\}

\begin{tblr}{width=\textwidth,colspec={|Q[l,m,3.2cm]|X[l,m]|},hlines,vlines,hspan=minimal,row{1-3}={bg=headgray},rowsep=1.5pt,colsep=3pt}
Jadwal Pelaksanaan: & Minggu 1--17 sesuai RPS. Setiap evaluasi memiliki RTM dan rubrik multi-kriteria. \\
Lain-Lain Yang Diperlukan: & LMS, repositori/kode sumber Python, perangkat lunak sesuai mata kuliah, dan perangkat presentasi. \\
Pustaka: & Mengacu pustaka utama dan pendukung pada bagian RPS. \\
\end{tblr}
